% sec-10: Document 10, parking and participant transportation standards.
\docpart{Parking and Transport}{PDAC LLM Oncology Clinical Trial Site Parking
and Patient Transportation Standards}

\section{Scope and Applicability}

This document governs parking supply, the accessible route, arrival and
departure geometry, and the transportation a participant needs after a
pancreaticoduodenectomy. It applies to the leasehold described in the front
matter and to the shared structure serving it.

Where the City of San Diego parking minimum for a medical office use and the
requirement derived here differ, the higher governs. That rule is stated once,
here, and is not repeated in document 05 \cite{sdmc}.

\section{Stall Count, Mix, and Dimensions}

The participant stall count is derived from the visit schedule rather than from
a floor area ratio. A Phase 1 site treating up to eighteen participants has a
peak concurrency that can be counted, and counting it is more defensible before
a hearing officer than applying a general medical office ratio to 12,400 square
feet.

\begin{mvtable}
\begin{tabularx}{\textwidth}{>{\raggedright\arraybackslash}p{3.2cm} >{\raggedright\arraybackslash}p{1.5cm} >{\raggedright\arraybackslash}p{2.6cm} >{\raggedright\arraybackslash}X}
\headrow \thead{Stall class} & \thead{Count} & \thead{Dimension} & \thead{Basis for the count}\\
\midrule
Participant & 22 & 8.5 by 18 feet & Peak concurrency of 6 participants, each with up to 2 accompanying vehicles, plus 4 for screening visits \\
Staff & 12 & 8.5 by 18 feet & 11 roles, of which at most 9 are on site at once, plus 3 for monitors and vendors \\
Accessible & 4 & 9 by 18 feet, one van-accessible at 11 feet & Above the code minimum for 46 stalls; an oncology population justifies the surplus \\
Electric vehicle & 3 & 8.5 by 18 feet & Title 24 new construction and alteration provisions \cite{ccr24} \\
Drop-off, short term & 3 & 8.5 by 20 feet & 20-minute limit; adjacent to the canopy \\
Loading and service & 2 & 10 by 25 feet & Deliveries and waste pickup under document 09 §5 \\
\textbf{Total} & \textbf{46} & & Front matter site table \\
\end{tabularx}
\end{mvtable}
\tabcapl{tab:parking}{The stall schedule. The count column sums to the 46 stalls
fixed in the front matter site table, and the basis column shows the arithmetic
rather than citing a ratio.}

\section{Accessible Parking and Path of Travel}

\begin{enumerate}
\item Four accessible stalls are provided, one of which is van-accessible, which
  exceeds the minimum required for a 46-stall facility under the 2010 ADA
  Standards for Accessible Design and Title 24 \cite{adastandards,ccr24}.

\item The accessible route runs from the accessible stalls to the canopy, into
  the lobby, and to the clinical zone entrance, with a running slope not
  exceeding 1:20 and a cross slope not exceeding 1:48, a minimum clear width of
  48 inches, and no step or lip exceeding one-half inch.

\item The route is signed at each decision point, and the signage is the
  wayfinding signage required by document 09 §6. There is one route, described
  once, in two documents.

\item Every door on the accessible route has a maximum opening force of 5 pounds
  and a clear width of not less than 32 inches.
\end{enumerate}

\section{Drop-Off, Ambulance, and Loading Geometry}

\begin{mvtable}
\begin{tabularx}{\textwidth}{>{\raggedright\arraybackslash}p{3.4cm} >{\raggedright\arraybackslash}p{3.0cm} >{\raggedright\arraybackslash}X}
\headrow \thead{Element} & \thead{Requirement} & \thead{Reason}\\
\midrule
Canopy vertical clearance & 14 feet 6 inches & A Type III ambulance with a roof-mounted air conditioning unit clears at 12 feet 6 inches; the margin covers a raised gurney \\
Ambulance approach & Separate from the participant drop-off lane & A transfer must not be blocked by a waiting vehicle \\
Turning radius, inside & 28 feet & Ambulance and the largest delivery vehicle \\
Drop-off lane length & 60 feet, two vehicles & Peak concurrency of 6 participants staggered at 30-minute intervals \\
Loading dock & At grade, with a levelling plate & Robot crate delivery during move-in; regulated waste pickup thereafter \\
Exterior queuing & Prohibited & Document 05 §5; all waiting is interior \\
\end{tabularx}
\end{mvtable}
\tabcapl{tab:geometry}{Six geometry requirements. The ambulance approach is
separated from the drop-off lane for one reason: the evacuation route in
document 11 §2 uses the same approach, and a route that can be blocked by a
parked car is not a route.}

\section{Post-Procedure Transportation}

A participant discharged after a pancreaticoduodenectomy does not drive. The
requirement is absolute, and the site does not treat it as advice.

\begin{enumerate}
\item Before a procedure is scheduled, the research nurse and participant
  navigator confirms and records the participant's transportation arrangement
  for the day of discharge, naming the driver or the service.

\item The confirmation is recorded in the visit record. A procedure is not
  scheduled without it.

\item On the day of discharge, the arrangement is re-confirmed before the
  discharge order is written.

\item Where the arrangement fails on the day, the site provides a chartered
  medical transport at its own cost, funded from the participant cost line in
  document 15 §1. The alternative, discharging a participant into a rideshare
  after a major abdominal resection, is not available.

\item Where a participant has no arrangement available for any scheduled visit,
  the navigator arranges transport in advance rather than at the door, and the
  cost is charged to the same line.
\end{enumerate}

The participant-facing framing here follows the patient advocacy work: a right
that depends on a participant arranging something difficult, unaided, is a right
in name \cite{patientadvocacy}.

\section{Transportation Demand Management and Monitoring}

\begin{mvtable}
\begin{tabularx}{\textwidth}{>{\raggedright\arraybackslash}p{4.0cm} >{\raggedright\arraybackslash}p{2.8cm} >{\raggedright\arraybackslash}X}
\headrow \thead{Measure} & \thead{Target} & \thead{Evidence}\\
\midrule
Staggered appointment scheduling & No more than 6 participant vehicles on site at once & The appointment book, sampled monthly \\
Staff transit subsidy & Offered to all 11 roles & Payroll record \\
Secure bicycle parking & 6 spaces & As-built drawing \\
Electric vehicle charging & 3 stalls, shared & Utilization log \\
Annual monitoring report & Filed with the City & The report, on the interval set by document 05 §6 \\
\end{tabularx}
\end{mvtable}
\tabcapl{tab:tdm}{Five demand management measures and the evidence for each. The
reporting interval is the interval in document 05 §6, so the City receives one
report rather than two on different schedules.}
